\chapter{Technique mentale}\label{ch:iv1:mental-technique}

\begin{ivtrains}
Les méthodes qui sous-tendent le calcul rapide~: pourquoi un raccourci fonctionne, quand il échoue et comment vérifier son résultat. Les recruteurs interrogent directement sur ces méthodes, et ils les attendent aussi au fil de questions plus longues. Le chapitre commence par les produits et les pourcentages, passe aux intérêts composés, aux logarithmes et aux racines, et se termine par des estimations qui combinent ces outils. Sans calculatrice, en énonçant chaque étape à voix haute.
\end{ivtrains}

\begin{question}[$\star$ \iqroles{QT, Trader, Analyste} \iqfirm{toutes} \iqlevel{junior}]\label{iq:iv1:mental-technique:1}
De tête, combien font 97 fois 103~? Qu'est-ce qui rend ce calcul rapide, et dans quels cas la même idée fonctionne-t-elle~?
\end{question}

\begin{question}[$\star$ \iqroles{QT, Trader} \iqfirm{teneur de marché, banque} \iqlevel{junior}]\label{iq:iv1:mental-technique:2}
Combien vaut 85 au carré~? Il existe une règle rapide pour élever au carré tout nombre qui se termine par 5~: énoncez-la et démontrez-la.
\end{question}

\begin{question}[$\star$ \iqroles{QT, Trader} \iqfirm{teneur de marché} \iqlevel{junior}]\label{iq:iv1:mental-technique:3}
Les nombres 97 et 96 sont tous deux juste en dessous de 100. Montrez comment les multiplier en n'utilisant que l'écart de chacun d'eux à 100, et expliquez pourquoi votre méthode fonctionne.
\end{question}

\begin{question}[$\star$ \iqroles{QT, Trader, Analyste} \iqfirm{toutes} \iqlevel{junior}]\label{iq:iv1:mental-technique:4}
De tête, combien font 368 fois 125~? Quelle astuce générale utilisez-vous, et quels autres multiplicateurs couvre-t-elle~?
\end{question}

\begin{question}[$\star$ \iqroles{QT, Trader} \iqfirm{teneur de marché, banque} \iqlevel{junior}]\label{iq:iv1:mental-technique:5}
Il existe une astuce pour multiplier par 11 en additionnant les chiffres voisins. Appliquez-la à 486 fois 11, en tenant compte des retenues, et expliquez pourquoi elle fonctionne.
\end{question}

\begin{question}[$\star$ \iqroles{QT, Trader, Analyste, Risque} \iqfirm{toutes} \iqlevel{junior}]\label{iq:iv1:mental-technique:6}
Combien font 16~\% de 25~? Pourquoi peut-on permuter les deux nombres dans une question de ce type~? Appliquez la même idée à 18~\% de 50, puis à 4~\% de 175.
\end{question}

\begin{question}[$\star$ \iqroles{QT, Trader, Analyste, Risque} \iqfirm{toutes} \iqlevel{junior}]\label{iq:iv1:mental-technique:7}
Un prix de 384 inclut une taxe de 20~\% calculée sur le prix hors taxe. Quelle part des 384 représente la taxe~? Donnez la méthode la plus rapide.
\end{question}

\begin{question}[$\star$ \iqroles{QT, Trader, Analyste, Risque} \iqfirm{toutes} \iqlevel{junior}]\label{iq:iv1:mental-technique:8}
Un placement perd 20~\%. Quelle hausse en pourcentage le ramène à son niveau initial~? Donnez la formule générale pour une baisse de taille quelconque, et appliquez-la à une baisse de 37.5~\%.
\end{question}

\begin{question}[$\star$ \iqroles{QT, Trader, Analyste} \iqfirm{toutes} \iqlevel{junior}]\label{iq:iv1:mental-technique:9}
Lequel est le plus grand, $7/9$ ou $11/14$, et de combien~? Répondez sans transformer l'un ou l'autre en décimal.
\end{question}

\begin{question}[$\star$ \iqroles{QT, Trader} \iqfirm{teneur de marché} \iqlevel{junior}]\label{iq:iv1:mental-technique:10}
Vous savez que $1/7 = 0.142857142857\dots$, avec ces six chiffres qui se répètent. Comment écrivez-vous $3/7$ et $5/7$ sous forme décimale sans diviser, et pourquoi votre méthode fonctionne-t-elle~?
\end{question}

\begin{question}[$\star$ \iqroles{QT, Trader, Analyste} \iqfirm{toutes} \iqlevel{junior}]\label{iq:iv1:mental-technique:11}
Vous achetez 300 actions à 80 et 500 actions à 88. Quel est votre prix moyen~? Trouvez-le sans calculer l'un ou l'autre des coûts totaux, et dites quelle règle vous utilisez.
\end{question}

\begin{question}[$\star$ \iqroles{QT, Trader, Analyste, Risque} \iqfirm{toutes} \iqlevel{junior}]\label{iq:iv1:mental-technique:12}
Selon la règle de 72, combien de temps faut-il à une somme pour doubler à 8~\% par an, avec une capitalisation annuelle~? Quel est l'écart entre ce résultat et la réponse exacte, et pourquoi une règle de ce genre fonctionne-t-elle seulement~?
\end{question}

\begin{question}[$\star$ \iqroles{QT, Trader, Risque} \iqfirm{toutes} \iqlevel{junior}]\label{iq:iv1:mental-technique:13}
Une action a une volatilité annuelle de 32~\%, c'est-à-dire que son rendement annuel a un écart-type de 32~\%. Il y a 252 jours de cotation dans une année, et les rendements journaliers sont indépendants. Quel est l'écart-type du rendement d'une journée~? Par quel nombre rapide divisez-vous~?
\end{question}

\begin{question}[$\star$ \iqroles{QR, QT, Trader} \iqfirm{toutes} \iqlevel{junior}]\label{iq:iv1:mental-technique:14}
On vous indique seulement que $\log_{10} 2 \approx 0.301$ et $\log_{10} 3 \approx 0.477$. Calculez $\log_{10} 72$ et $\log_{10} 0.0072$ sans calculatrice.
\end{question}

\begin{question}[$\star$ \iqroles{QT, QD, Dev, LLE} \iqfirm{toutes} \iqlevel{junior}]\label{iq:iv1:mental-technique:15}
Sans calculatrice, à peu près combien vaut $2^{30}$, et combien de doublements faut-il pour passer de 1 à mille milliards~? De combien votre première estimation s'écarte-t-elle, et quel fait utilisez-vous~?
\end{question}

\begin{question}[$\star$ \iqroles{QR, QT, Trader} \iqfirm{toutes} \iqlevel{junior}]\label{iq:iv1:mental-technique:16}
Estimez $\sqrt{50}$ à trois décimales près, en une seule étape de calcul mental. Votre estimation est-elle trop haute ou trop basse, et pourquoi se situe-t-elle toujours de ce côté~?
\end{question}

\begin{question}[$\star$ \iqroles{QT, Trader, Analyste, Risque} \iqfirm{banque, hedge fund} \iqlevel{junior}]\label{iq:iv1:mental-technique:17}
Un point de base vaut un centième d'un pour cent. Combien valent 10 points de base de 3.5~millions de dollars~? Un prêt de 250~millions de dollars verse un supplément d'un point de base par an~: combien cela représente-t-il en dollars sur trois mois~? Donnez la méthode rapide.
\end{question}

\begin{question}[$\star$ \iqroles{QT, Trader, Risque} \iqfirm{banque, hedge fund} \iqlevel{junior}]\label{iq:iv1:mental-technique:18}
Le cours d'une obligation est coté en trente-deuxièmes de point, de sorte que 99-17 signifie 99 et $17/32$. Combien vaut 99-17 sous forme décimale, et comment l'obtenez-vous sans diviser par 32~?
\end{question}

\begin{question}[$\star\star$ \iqroles{QR, QT, Trader} \iqfirm{toutes} \iqlevel{junior}]\label{iq:iv1:mental-technique:19}
Sans calculer entièrement l'un ou l'autre produit, lequel est le plus grand, $37 \times 43$ ou $38 \times 42$, et de combien exactement~? Quelle est la règle générale~?
\end{question}

\begin{question}[$\star\star$ \iqroles{QT, Trader, Risque} \iqfirm{toutes} \iqlevel{junior}]\label{iq:iv1:mental-technique:20}
Calculez $1.02 \times 0.97$ à quatre décimales près, de tête. Quelle partie du calcul pouvez-vous négliger lorsque vous n'avez besoin que d'une réponse approchée, et quelle est l'ampleur de l'erreur~?
\end{question}

\begin{question}[$\star\star$ \iqroles{QR, QT, Trader, QD, Dev} \iqfirm{toutes} \iqlevel{junior}]\label{iq:iv1:mental-technique:21}
Quelqu'un vous affirme que $3847 \times 26 = 100\,122$. Montrez que c'est faux sans effectuer la multiplication. Quelles vérifications rapides cette égalité passe-t-elle, et lesquelles échoue-t-elle~?
\end{question}

\begin{question}[$\star\star$ \iqroles{QT, Trader} \iqfirm{teneur de marché} \iqlevel{junior}]\label{iq:iv1:mental-technique:22}
Combien vaut $73 \times 68$~: 4964, 4954, 5064 ou 4864~? Les quatre se terminent par le bon dernier chiffre et ont à peu près le bon ordre de grandeur. Comment choisissez-vous la bonne réponse sans multiplier~?
\end{question}

\begin{question}[$\star\star$ \iqroles{QR, QT, QD} \iqfirm{toutes} \iqlevel{junior}]\label{iq:iv1:mental-technique:23}
Quelqu'un écrit $71 \times 88 = 6428$. L'addition des chiffres et la comparaison modulo 9 ne détectent pas cette erreur. Pourquoi cela, et quelle vérification rapide permet de la repérer~?
\end{question}

\begin{question}[$\star\star$ \iqroles{QT, Trader, Analyste, Risque} \iqfirm{toutes} \iqlevel{junior}]\label{iq:iv1:mental-technique:24}
Un prix augmente de 20~\%, puis baisse de 15~\%. Quelle est la variation nette, et l'ordre a-t-il de l'importance~? Donnez une règle pour combiner de tête deux variations en pourcentage.
\end{question}

\begin{question}[$\star\star$ \iqroles{QR, QT, Trader} \iqfirm{toutes} \iqlevel{junior}]\label{iq:iv1:mental-technique:25}
Estimez $1/0.97$ et $1/1.03$ à quatre décimales près, sans diviser. Quelle approximation utilisez-vous, et à partir de quand cesse-t-elle de fonctionner~?
\end{question}

\begin{question}[$\star\star$ \iqroles{QR, QT, Trader, Risque} \iqfirm{toutes} \iqlevel{junior}]\label{iq:iv1:mental-technique:26}
Estimez $1.05^{10}$ à deux décimales près, de tête, de deux manières différentes.
\end{question}

\begin{question}[$\star\star$ \iqroles{QR, QT, Trader, Risque} \iqfirm{teneur de marché, hedge fund} \iqlevel{junior}]\label{iq:iv1:mental-technique:27}
Estimez $0.99^{100}$ à trois décimales près, sans calculatrice. Vous pouvez utiliser $1/e \approx 0.368$. Pourquoi $1/e$ seul ne suffit-il pas~?
\end{question}

\begin{question}[$\star\star$ \iqroles{QR, QT, Trader} \iqfirm{toutes} \iqlevel{junior}]\label{iq:iv1:mental-technique:28}
La règle du 72 affirme qu'un capital double en environ $72/r$ années à un taux de $r$ pour cent par an, avec capitalisation annuelle. Quelle est son erreur pour $r = 1$ et pour $r = 36$~? Au voisinage de quel taux est-elle exacte~? Quelle correction simple permet de la corriger pour tous ces taux~?
\end{question}

\begin{question}[$\star\star$ \iqroles{QR, QT, Trader, Risque} \iqfirm{teneur de marché, hedge fund} \iqlevel{junior}]\label{iq:iv1:mental-technique:29}
Votre compte augmente de 10 points de base chaque jour pendant 252 jours, un point de base valant 0.01~\%. Estimez de tête le rendement total à un dixième de point de pourcentage près. Quelle est la réponse simple, et pourquoi est-elle trop faible~?
\end{question}

\begin{question}[$\star\star$ \iqroles{QR, QT, Trader, Risque} \iqfirm{teneur de marché, hedge fund} \iqlevel{junior}]\label{iq:iv1:mental-technique:30}
Votre investissement augmente de 10~\% puis baisse de 10~\%, tour à tour, pendant 100 périodes. Quelle fraction de votre capital reste-t-il à la fin~? Le rendement moyen par période est nul, alors pourquoi avez-vous perdu de l'argent~? Énoncez la règle qui sous-tend la réponse.
\end{question}

\begin{question}[$\star\star$ \iqroles{QT, Trader, Risque} \iqfirm{toutes} \iqlevel{junior}]\label{iq:iv1:mental-technique:31}
Une position perd 3~\% de sa valeur restante chaque année. Environ combien de temps faut-il pour que la moitié ait disparu~? La règle du 72 donne-t-elle un temps trop long ou trop court ici, et en quoi cela diffère-t-il du cas du doublement au même taux~?
\end{question}

\begin{question}[$\star\star$ \iqroles{QR, QT} \iqfirm{toutes} \iqlevel{junior}]\label{iq:iv1:mental-technique:32}
On vous indique seulement que $\log_{10} 2 \approx 0.301$ et $\log_{10} 3 \approx 0.477$. Trouvez $\log_{10} 7$ à trois décimales près, sans calculatrice.
\end{question}

\begin{question}[$\star\star$ \iqroles{QR, QT, QD, Dev} \iqfirm{toutes} \iqlevel{junior}]\label{iq:iv1:mental-technique:33}
Combien de chiffres décimaux possède $2^{100}$, et quel est son premier chiffre~? Vous pouvez utiliser $\log_{10} 2 \approx 0.30103$.
\end{question}

\begin{question}[$\star\star$ \iqroles{QR, QT, Trader} \iqfirm{toutes} \iqlevel{junior}]\label{iq:iv1:mental-technique:34}
Vous savez que $\ln 2 \approx 0.693$ et $\ln 10 \approx 2.303$. Estimez $\ln 1.3$ à trois décimales près. Pourquoi la série de $\ln(1 + x)$ seule ne vous est-elle pas d'un grand secours ici~?
\end{question}

\begin{question}[$\star\star$ \iqroles{QR, QT, Trader} \iqfirm{toutes} \iqlevel{junior}]\label{iq:iv1:mental-technique:35}
Estimez la racine dixième de 2 et la racine cinquième de 3, chacune à trois décimales près, sans calculatrice. Vous pouvez utiliser $\ln 2 \approx 0.693$ et $\ln 3 \approx 1.099$.
\end{question}

\begin{question}[$\star\star$ \iqroles{QT, Trader, Analyste} \iqfirm{toutes} \iqlevel{junior}]\label{iq:iv1:mental-technique:36}
Estimez $47 \times 62 \times 81 \times 19$ à 1~\% près, sans écrire aucun produit intermédiaire. Comment corrigez-vous l'arrondi~?
\end{question}

\begin{question}[$\star\star$ \iqroles{QR, QT, Trader, DS} \iqfirm{toutes} \iqlevel{junior}]\label{iq:iv1:mental-technique:37}
Vous lancez une pièce équilibrée 20 fois. Estimez de tête la probabilité d'obtenir exactement 10 piles, sans calculer de coefficient binomial. Estimez ensuite la probabilité d'obtenir exactement 12 piles.
\end{question}

\begin{question}[$\star\star\star$ \iqroles{QR, QT} \iqfirm{teneur de marché, hedge fund} \iqlevel{junior, confirmé}]\label{iq:iv1:mental-technique:38}
De combien de façons peut-on choisir 25 éléments parmi 50~? Donnez deux chiffres significatifs, sans calculatrice. Faites de même pour choisir 30 éléments parmi 100.
\end{question}

\begin{question}[$\star\star\star$ \iqroles{QR, QT, QD, Dev} \iqfirm{toutes} \iqlevel{junior, confirmé}]\label{iq:iv1:mental-technique:39}
Combien de bits faut-il pour enregistrer l'ordre d'un jeu de 52 cartes mélangé, avec le moins de bits possible~? Estimez-le sans calculatrice. Combien de chiffres décimaux possède $52!$~?
\end{question}

\begin{question}[$\star\star\star$ \iqroles{QR, QT, Trader, Analyste, Risque} \iqfirm{banque, hedge fund} \iqlevel{junior, confirmé}]\label{iq:iv1:mental-technique:40}
Vous recevez 1~\$ à la fin de chaque année pendant 30 ans, et le taux d'actualisation est de 7~\% par an. Estimez la valeur actuelle à une décimale près, sans calculatrice. Calculez ensuite le prix d'une obligation de 30 ans qui verse un coupon annuel de 5~\%, au même taux.
\end{question}
